% TOP_PROBLEM: {"id":30007595,"problem_number":"LOCAL-30007595","title":"Belief aggregation","category_id":21,"category":{"id":21,"name":"economics","display_name":"Economics","description":"Open economic theory statements, published research agendas, and proposed scoped specifications; question provenance and historical priority are qualified per record.","slug":"economics","order_index":21,"created_at":"2026-10-08T00:00:00Z"},"status":"research_question","external_url":null,"legacy_ids":[],"published":false,"statement":"In the explicitly specified setting: How do heterogeneous subjective beliefs and perceived risks aggregate into equilibrium prices and trading? The mathematical statement above fixes the target, assumptions and scope of this catalogue formulation.","economics":{"original_id":84,"field":"Asset pricing and financial markets","kind":"theoretical","formalization_basis":"adapted_research_question","source_status":"explicit_open","priority_status":"earliest_found","priority_note":"This is the earliest explicit statement verified in this audit; first priority for the full catalogue formulation is not established.","earliest_verified_reference":{"authors":["Klaus Adam","Stefan Nagel"],"title":"Expectations Data in Asset Pricing","year":2022,"url":"https://bpb-us-w2.wpmucdn.com/voices.uchicago.edu/dist/f/575/files/2022/01/ExpectationsSurvey_13.pdf","doi":"10.3386/w29977","locator":"Section 6, Future Research Directions, pp.35–36; abstract","evidence_summary":"Calls for belief aggregation and subjective-risk evidence."},"current_reference":{"authors":["Klaus Adam","Stefan Nagel"],"title":"Expectations Data in Asset Pricing","year":2022,"url":"https://bpb-us-w2.wpmucdn.com/voices.uchicago.edu/dist/f/575/files/2022/01/ExpectationsSurvey_13.pdf","doi":"10.3386/w29977","locator":"Section 6, Future Research Directions, pp.35–36; abstract","evidence_summary":"Calls for belief aggregation and subjective-risk evidence."},"formalization_note":"Adapts the source’s explicit belief-aggregation agenda. The survey measurement and equilibrium-map specifications are proposed, and require model choices before implementation."},"statusQualification":"Published question or research agenda verified at the cited locator. The mathematical specification is adapted for this catalogue and includes additional assumptions; source publication does not establish first-ever priority or certify this exact model remains unsolved. Adapts the source’s explicit belief-aggregation agenda. The survey measurement and equilibrium-map specifications are proposed, and require model choices before implementation.","statusReviewedAt":"2026-10-08"}
% FIRST_POSTED: 2026-10-08
% ENTRY_KIND: statement_only
% !TeX program = lualatex
\documentclass[11pt]{article}
\usepackage{fontspec}
\usepackage[a4paper,margin=25mm]{geometry}
\usepackage{amsmath,amssymb,mathtools}
\usepackage{xurl}
\usepackage[hidelinks]{hyperref}
\newcommand{\sref}[2]{\hyperlink{source:#1:#2}{[S#2]}}
\setlength{\emergencystretch}{3em}
\begin{document}
% problemId:30007595
\section*{Belief aggregation}\label{30007595}
\subsection{Definitions and mathematical statement}
Consider a finite-horizon market with agents \(i=1,\ldots,N\), asset dividends \(D_t\), prices \(P_t\), wealth \(w_{it}\), utility functions and explicitly specified portfolio constraints. Agent \(i\) evaluates uncertainty using a subjective conditional law \(\mathbb P_{it}\), which need not coincide with the objective data-generating law. Survey responses \(b_{it}\) are noisy measurements of these laws under a stated measurement model. An equilibrium price process is generated by optimal demands \(q_{it}\) and clearing \(\sum_iq_{it}=Q_t\), where \(Q_t\) is exogenous supply. For an unconstrained marginal agent, pricing satisfies
\[P_t=E_{\mathbb P_{it}}[M_{i,t+1}(P_{t+1}+D_{t+1})\mid\mathcal F_{it}],\]
with \(M_{i,t+1}\) the agent's marginal-utility discount factor and \(\mathcal F_{it}\) their information. The task is to characterize an empirically estimable aggregation map from the joint distribution of beliefs, risk perceptions, wealth and constraints to prices and volume. Determine conditions under which survey observations identify this map or its relevant derivatives. Explicitly allow the identity of marginal investors to change with prices and constraints. Quantify sensitivity to unobserved nonparticipants, reporting error and selective survey response. Matching average beliefs is insufficient if price-setting agents have different beliefs or risk tolerance.

\par\textbf{Formulation and scope.} Adapts the source’s explicit belief-aggregation agenda. The survey measurement and equilibrium-map specifications are proposed, and require model choices before implementation.

\par\textbf{Open-question provenance.} An explicit question or research agenda was verified at the source locator. This does not by itself certify that every later version remains unresolved \sref{30007595}{1}.

\par\textbf{Historical priority.} This is the earliest explicit statement verified in this audit; first priority for the full catalogue formulation is not established.

\subsection{Short English statement}
In the explicitly specified setting: How do heterogeneous subjective beliefs and perceived risks aggregate into equilibrium prices and trading? The mathematical statement above fixes the target, assumptions and scope of this catalogue formulation.

\subsection{Sources}
\begin{itemize}\raggedright
\item[\textnormal{[S1]}]\hypertarget{source:30007595:1}{} Klaus Adam, Stefan Nagel. 2022. Expectations Data in Asset Pricing. Section 6, Future Research Directions, pp.35–36; abstract.
\url{https://bpb-us-w2.wpmucdn.com/voices.uchicago.edu/dist/f/575/files/2022/01/ExpectationsSurvey_13.pdf}.
DOI: 10.3386/w29977.
{\small\textit{Earliest verified question/agenda reference; historical priority qualified below. Calls for belief aggregation and subjective-risk evidence.}}
\end{itemize}
\end{document}
